% Ad Atticum 9.13 --- headnote
% ad-atticum-09-13, 24 March 49 BC, Formian villa

\textit{Cicero to Atticus, written from the Formian villa
on the ninth day before the Kalends of April 49 BC (the
manuscript dateline: \textit{Scr.\ in Formiano ix K.\ Apr.\
a.\ 705 (49)}). Four days after the panic of 9.12, the
rafts-at-Brundisium story has been overtaken by a fresh
dispatch from Dolabella, dated 13 March from Brundisium:
this is Caesar's \textit{[Greek: euh\=emer\'ia]} ---
his ``fine day'' --- and the report is that Pompey is
in flight and will sail with the first wind. Cicero
opens with the Homeric tag \textit{[Greek: ouk \'est'
\'etumos l\'ogos]}, ``the tale is not true'' (Stesichorus's
formula known from Plato), and reconstructs the
chronology: the day before the Quinquatria there was
a great storm-wind, and that, he supposes, is what
Pompey rode out. The siege-wall and the rafts were a
ghost; Pompey has been at sea since the 17th.}

\textit{Section 3 is the long, exposed self-defence of
Cicero's relationship to Caesar --- partly to Atticus,
partly to himself: ``I have always cried up his services
to me, the more lest he think I was remembering the
older grievances''; ``he gave me no help when he could,
but afterwards was a friend, even very much so'';
``therefore I in turn am his friend''; and then the
clean accounting: ``I do not now know in what way I
could help him; nor, if I could, when he is preparing
so destructive a war, should I think I ought to.'' What
he dreads is not Caesar's \textit{[Greek: go\=ete\'ia]},
his wizardry, but his \textit{[Greek: peithan\'ank\=e]} ---
persuasion-that-compels; and the line is sealed with a
direct quotation of Plato's \textit{Seventh Letter}:
``the requests of tyrants are mixed with compulsion.''
Section 4 carries a corrupt clause, marked here with
daggers as in the standard text. The closing sections
record Lentulus Spinther still at Puteoli, \textit{[Greek:
ad\=emon\^on]} (at his wits' end) and dreading another
Corfinium; Balbus's letter, with its last paragraph in
which the financier is now ``in agony''; and the verdict
from Dolabella that the war will be \textit{merum bellum},
unmixed. Cicero closes by settling into the desperation
itself: ``let us stay in that same wretched and
despairing mind, since nothing can be more wretched than
this.''}
